\subsubsection*{注 5.1.}
条件

\begin{equation}\tag{5.1.1}
jd<n+2
\end{equation}

具有上同调解释. 设 $V$ 为次数 $d$、维数 $n$ 的光滑超曲面.
如 XI, 2.1 中所做, 对 $a+b=n$, 置

\begin{equation}\tag{5.1.2}
h_0^{a,b}
=\dim H^b(V,\Omega_V^a)-\delta_{a,b}
=\dim E^{a,b}(V),
\end{equation}

其中 $E^{a,b}(V)$ 为消灭 Hodge 上同调的双次数 $(a,b)$ 分量 (比较 XVIII, 5.4.4)\EditorialFootnote{原书第430页及英文底本作 XVII, 5.4.4; 本卷对消灭上同调空间的定义见 XVIII, 5.4.4.}

\begin{equation}\tag{5.1.3}
E^{a,b}(V)
=\operatorname{Coker}\!\left(
H^b(\mathbb P^{n+1},\Omega_{\mathbb P^{n+1}}^a)
\longrightarrow H^b(V,\Omega_V^a)\right).
\end{equation}

则由 XI, 2.8,

\begin{equation}\tag{5.1.4}
jd<n+2
\quad\Longleftrightarrow\quad
h_0^{a,n-a}=0\quad\text{当 }a<j\text{ 或 }a>n-j.
\end{equation}

\begin{editorialnote}
原书第430页及英文底本写作上述双向断言. 引用 XI, 2.8 时须注意其中
$0\leq\ell<n$ 的范围; 置 $\ell=n-2j$ 后为 $1\leq j\leq n/2$.
不能将这里的逆向断言无条件延伸到该范围之外: 例如光滑平面二次曲线有 $n=1$, $d=2$,
其消灭 Hodge 上同调为零, 故取 $j=2$ 时右端仍成立, 而左端成为不成立的 $4<3$.
此处保留原书的双向表述, 同时明确这一适用范围问题.
\end{editorialnote}

换言之, 此条件说明 $V$ 的消灭 Hodge 上同调

\[
E_{\mathrm{DR}}^n(V)
=\operatorname{Coker}\!\left(
H_{\mathrm{Hdg}}^n(\mathbb P^{n+1})
\longrightarrow H_{\mathrm{Hdg}}^n(V)\right),
\]

的 Hodge 层次 (XI, 2.7) 至多为 $n-2j$, 或等价地其余层次至少为 $j$.
另一方面, 结论也说明, Frobenius $\mathfrak F_q$ 在消灭 $\ell$ 进上同调
($\ell\ne p$) 上作用的特征值是被 $q^j$ 整除的代数整数 (XXI, 5.2.2).\EditorialFootnote{原书第431页及英文底本沿用“XXI, A.2”的旧引文; 此处所需的上同调整性定理在本卷编号为 XXI, 5.2.2.}
因此 Ax 的结果 5.0 支持下述猜想, 它由动机层次的思想所启发,
并联系着经典的 Hodge 和 Tate 猜想~[5].

\subsubsection*{猜想 5.2.}
设 $f:X\to S$ 为固有光滑态射, $S$ 为整概形, 在 $\operatorname{Spec}(\mathbb Z)$ 上有限型,
且支配 $\operatorname{Spec}(\mathbb Z)$; 设 $i,j$ 为整数.
假定 $f$ 的泛纤维的 $H^i$ 的 Hodge 余层次至少为 $j$.
则对每个点 $s\in S$, $H^i(X_s,\mathbb Q_\ell)$ 的余层次 (XX, 2.0.6) 至少为 $j$.
因此 (XXI, 2.2),\EditorialFootnote{原书第431页及英文底本作 XX, 2.2; 此处所用的余层次与 Frobenius 整除性的推论为 XXI, 2.2. XX, 2.2 则说明代数闭域扩张下滤过的不变性.} 若 $s$ 是 $S$ 的闭点, Frobenius 在 $H^i(X_s,\mathbb Q_\ell)$
上作用的特征值是被 $q^j$ 整除的代数整数, 其中 $q=\operatorname{card}(k(s))$.

\section*{6. 同余公式证明的结束}

考虑 (4.6.6) 及 (4.6.3)--(4.6.4), 公式 (4.5.33) 给出

\begin{equation}\tag{6.0.1}
Z(t,V/\mathbb F_q)(1-t)
\equiv\Delta_S(t)^{(-1)^{n+1}}
\pmod{1+qt\mathcal O[[t]]}.
\end{equation}

现在计算 $\Delta_S(t)$ 模 $\pi\mathcal O[[t]]$ 的结果. 如 (4.6.11) 中那样定义

\begin{equation}\tag{6.0.2}
W_S=R_S\otimes_{\mathcal O}\mathbb F_q.
\end{equation}

\medskip
\noindent\emph{底本注.}
印本引用 (4.5.24), 但该同余式由末尾的 zeta 恒等式 (4.5.33) 得出.
它还用 $R_B$ 定义 $W_S$, 而全节及随后每个公式均讨论 $R_S$.
校订法文层和英文层采用 (4.5.33) 及 $R_S$; 照录法文层保留印本形式.
